% !TEX root = chapter-standalone.tex

\section{Broad ToDos and comments}

(This section created in June 2026)


\subsection{Meta comments}

\begin{enumerate}
\item In general, we should watch out that on any given page or short sequence of pages there are not too many main messages being communicated at once. The main point should be clear, and larger distractions from the main line of storytelling and explanation should be avoided. 
\item Generally it would be good to illustrate various of the ``simple mathematical examples'' with further, more concrete and applied examples -- in order to spell out the more concrete implications and possible applied uses of the mathematical examples. These additional concrete examples should be small enough that readers can easily chew on them, but insightful enough to be relevant. Even if the applied examples are relatively simple, this step from generic math to context specific use is often a gap that is valuable (and often nontrivial) to bridge). 
\item Overall, the book should not be thought of as a reference text and become to rigid in any attempt to be systematic or exhaustive; on the other hand, we want to avoid obscuring inherent systematic structures in category theory (and generally in our covered topics) that are helpful or even essential for readers to see and understand. 
\end{enumerate}


\subsection{Broad ToDos}

\todotext{Put the discussion of diagrams of sets and functions (that encode for example the equations in the definitions of semigroup, monoid, and group) into its own section}

\todotext{Remove quantifiers where type-inference makes them technically superfluous.}

\todotext{Make a pass through whole book and put into 'devel' any material that looks like it is very 'bare bones' and does not look like a semi-readable book passage.}

\todotext{draft material on quaternions (and octonions?) to consider incorporating in the book}
